% SHARD-ID: RC-02B-DEFINITIONS-ARITHMETIC
% SHARD-TITLE: Definitions, Continued: Riemann Channel, Weil--LPS, Curves
% SHARD-SUMMARY: Continues the single-occurrence definitions: the Riemann channel and Bost--Connes objects, the Weil--LPS construction, Artin--Schreier curves and Deligne's vocabulary.
% SHARD-SUMMARY: Split from the definitions shard only for length; every row of db/definitions.tsv still appears exactly once.
% SHARD-KEYWORDS: definitions, Riemann channel, scattering symbol, Bost-Connes, Weil representation, LPS, Artin-Schreier, Dwork, pure weight
% SHARD-NOTES: none
% SHARD-SCRIPTS: none

\section{Definitions, Continued}
\label{sec:definitions-b}

Continuation of \cref{sec:definitions}; split for length only.

\subsection{The Riemann channel}

\begin{definition}[Scattering symbol]
\label{def:scattering-symbol}
With $\cxi(s) = \frac12 s(s-1)\pi^{-s/2}\Gamma(s/2)\zeta(s)$ and $\tauM$ the
Mellin variable on the dilation group, the \emph{scattering symbol} of the
modular surface is $\Ssym(\tauM) = \cxi(1-2i\tauM)/\cxi(1+2i\tauM)$.
\end{definition}

\begin{definition}[Inner function, model space, compressed semigroup]
\label{def:inner-model-space}
A function analytic and bounded by $1$ in the lower half plane with
unimodular boundary values is \emph{inner}. For inner $\Ssym$ the
\emph{model space} is $\KS = H^2\ominus\Ssym H^2$ and the \emph{compressed
semigroup} is $\Zt = P_{\KS}\,e^{i\tsg\tauM}|_{\KS}$, $\tsg>0$.
\end{definition}

\begin{definition}[Riemann channel]
\label{def:riemann-channel}
The \emph{Riemann channel} is the completely positive semigroup
$\dm\mapsto\Zt^{*}\dm\Zt$ on $B(\KS)$, whose single Kraus operator is the
compressed semigroup of the scattering symbol. Its modes are the zeros of
$\zeta$ and its distributional trace is the prime measure; the name is also
used loosely for the whole side-A/side-B package of the Riemann zeta.
\end{definition}

\begin{definition}[Weil explicit formula]
\label{def:explicit-formula}
For a test window $\win$ with transform $g$,
$\sum_{\gamma}\win(\gamma) = \win(\tfrac i2)+\win(-\tfrac i2) - g(0)\log\pi
+ \frac1{2\pi}\int\win(r)\,\re\,\digam(\tfrac14+\tfrac{ir}2)\,dr
- 2\sum_n\frac{\vM(n)}{\sqrt n}\,g(\log n)$, the sum on the left over the
ordinates of the nontrivial zeros.
\end{definition}

\begin{definition}[Bost--Connes system, Lévy exponent]
\label{def:bost-connes}
The \emph{Bost--Connes system} is the $C^*$-dynamical system generated by the
phase operators $e(a/b)$ and the isometries $\bciso{n}$ on $\ltwo{\mathbb N}$,
with time evolution scaling $\bciso{n}$ by $n^{it}$; its KMS$_\invtemp$ states for
$\invtemp>1$ have partition function $\zeta(\invtemp)$ and $\invtemp = 1$ is its
phase transition. Its \emph{Lévy exponent} at real part $\sigs$ is
$\levy{\sigs}(\tauM) = \log\zeta(\sigs-i\tauM) - \log\zeta(\sigs)$, the
cumulant generating function of the zeta distribution, with jump measure
$\jump{\sigs}$.
\end{definition}

\begin{definition}[Matrix product operator over the prime chain]
\label{def:prime-chain-mpo}
Writing $\ltwo{\mathbb N} = \bigotimes_{\pp}$ (oscillator at $\pp$), an
operator is a \emph{matrix product operator over the prime chain} when it
is a contraction of one local tensor per prime with a bond space carried
between consecutive primes; for $e(a/b)$ the bond space is $\mathbb Z/b$ and
the local transfer map is $r\mapsto\pp^k r$.
\end{definition}

\begin{definition}[Hilbert--P\'olya operator]
\label{def:hilbert-polya}
On the subleading sector of the length flow write $\flowgen = -\tfrac12 +
iH$; $H$ is the \emph{Hilbert--P\'olya operator}, with eigenvalues $\gamma_n
- i(\sigs_n-\tfrac12)$. It is not the entanglement Hamiltonian of any state.
\end{definition}

\begin{definition}[Weil positivity]
\label{def:weil-positivity}
\emph{Weil's criterion}: the Riemann hypothesis holds iff
$\Fweil(\uz) = \sum_{\zero}e^{(\zero-1/2)\uz}$ is positive-definite on the
dilation group, i.e.\ iff its spectral measure $\sum_{\zero}\delta_{t_{\zero}}$
is a positive measure on the real line.
\end{definition}

\subsection{Weil--LPS channels}

\begin{definition}[LPS generators]
\label{def:lps-generators}
For a prime $\qs\equiv1\pmod4$, the \emph{LPS generators} are the $\qs+1$
integer quaternions $\quat = a_0+a_1i+a_2j+a_3k$ of norm $\qs$ with $a_0>0$
odd and $a_1,a_2,a_3$ even; they generate a $(\qs+1)$-regular tree, and their
images $\galpha$ in $\PSL(\Fp)$ (for $\qs$ a square mod $\pp$, after division
by $\sqrtq$) are the generators of the LPS Cayley graph.
\end{definition}

\begin{definition}[Weil representation]
\label{def:weil-representation}
The \emph{Weil representation} $\Weil$ of $\SL(\Fp)$ on $\ltwo{\Fp}$ is the
projective unitary representation characterised by $\Weil(g)\,\Hei\,
\Weil(g)^\dagger\propto T(g(u,v))$ for the Heisenberg--Weyl operators
$\Hei = \shift^u\clock^v$; it is generated by $\Four$, $\chirp$ and $\Dil$ and
splits into the blocks $\Wpm$ of dimensions $(\pp\pm1)/2$ on which $-1$ acts
as a scalar.
\end{definition}

\begin{definition}[Weil--LPS channel]
\label{def:weil-lps-channel}
For admissible $(\pp;\qs)$ the \emph{Weil--LPS channel} is
$\PhiWL{\qs}^{\pm}(\dm) = \frac1{\qs+1}\sum_{\quat}\Wpm(\galpha)\,\dm\,
\Wpm(\galpha)^\dagger$ on $M_{(\pp\pm1)/2}$: Harrow's construction with the
Weil representation and LPS generators.
\end{definition}

\begin{definition}[Hecke relation]
\label{def:hecke-relation}
The adjacency operators $\Hecke{\qs}$ of the LPS Cayley graph satisfy
$\Hecke{\qs}^2 - \qs\,1 = \Hecke{\qs^2}$ and $\Hecke{\qs_1}\Hecke{\qs_2} =
\Hecke{\qs_1\qs_2}$, the tree relations ``two steps without backtracking''
and multiplicativity; the Weil--LPS channels for different $\qs$ commute.
\end{definition}

\subsection{Curves}

\begin{definition}[Artin--Schreier curve, exponential sum]
\label{def:artin-schreier}
For $g\in\Fq[x]$ the \emph{Artin--Schreier curve} is $y^{\qs}-y = g(x)$; its
affine point count over $\Fqn{n}$ is $\Ncount{n} = \qs^n +
\sum_{a\in\Fq^\times}\expsum{n}(ag)$ with the \emph{exponential sum}
$\expsum{n}(g) = \sum_{x\in\Fqn{n}}\addch(\ftr\,g(x))$. The quadratic family
is $g(x) = \sum_{j=0}^{\asrange}a_jx^{1+\qs^j}$.
\end{definition}

\begin{definition}[Self-dual normal basis]
\label{def:self-dual-normal-basis}
A basis $\{\nbgen,\nbgen^{\qs},\dots,\nbgen^{\qs^{n-1}}\}$ of $\Fqn{n}$ is
\emph{normal}; Frobenius is then the cyclic shift of coordinates. It is
\emph{self-dual} when $\ftr(ab) = \sum_i a_ib_i$; such bases exist for $n$
odd.
\end{definition}

\begin{definition}[Dwork transfer operator]
\label{def:dwork-operator}
\emph{Dwork's operator} is $\dwork = \decim\circ M_F$ on $\pp$-adic power
series, multiplication by the splitting function $M_F$ followed by the decimation
$\decim: x^m\mapsto x^{m/\qs}$; $\expsum{n}^{\times} = (\qs^n-1)\Tr\dwork^n$ and
$\Lcurve$ is its Fredholm determinant.
\end{definition}

\begin{definition}[Pure of weight, big monodromy, slicing, the sign]
\label{def:pure-weight}
A coefficient system on a curve over $\Fq$ is \emph{pure of weight $\wt$}
when every eigenvalue of every holonomy of a prime cycle of length $\wl$ has
modulus $\qs^{\wt\wl/2}$; it has \emph{big monodromy} when the only
invariants of its even tensor powers are the pairwise contractions by the
given pairing. \emph{Slicing} writes the middle block of a variety as the
interesting block of a base curve with coefficients from hyperplane
sections. \emph{The sign} is the minus sign with which the interesting block
enters the point count, so that its eigenvalues are zeros, not poles, of
$\curvezeta$.
\end{definition}

\begin{definition}[Status vocabulary]
\label{def:status-vocabulary}
The tags \texttt{proved}, \texttt{sketched}, \texttt{numerical},
\texttt{cited}, \texttt{conjectured}, \texttt{open}, \texttt{refuted},
\texttt{assumed}, with the derived forms \texttt{unsupported} and
\texttt{-conditional}, as fixed in \cref{sec:conventions} of the frontmatter;
definitions are \texttt{stipulated} or \texttt{quoted}.
\end{definition}

\subsection{The continuous dictionary}

\begin{definition}[Geodesic flow on the unit tangent bundle]
\label{def:geodesic-flow}
For $\Gamma<\PSL(\mathbb R)$ discrete, torsion-free and cocompact, the unit
tangent bundle of $\Gamma\backslash\mathbb H$ is $\Gamma\backslash G$ with
$G = \PSL(\mathbb R)$, and the \emph{geodesic flow} is $\phi_t(\Gamma g) =
\Gamma g a_t$, $a_t = \operatorname{diag}(e^{t/2},e^{-t/2})$. Its generator
is $\gflow = \Hgen/2$ through the left-invariant fields $\Hgen,\Egen,\Wgen$ of
the basis $H,E,W$ of $\mathfrak{sl}_2$, and the Koopman operator is
$(e^{-t\gflow}f)(m) = f(\phi_{-t}m)$. Primitive closed geodesics $\gamma$
have lengths $\wl_\gamma$ and opposite orientations are not identified.
\end{definition}

\begin{definition}[Selberg zeta, Ruelle zeta]
\label{def:selberg-zeta}
$\Zsel(s) = \prod_\gamma\prod_{k\ge0}\big(1-e^{-(s+k)\wl_\gamma}\big)$, in
Marklof's words \texttt{\detokenize{Z(s) = \prod_{\gamma\in H_*} \prod_{m=0}^\infty \left( 1-\e^{-\ell_\gamma(s+m)} \right),}} (provenance \texttt{prov:marklof04-selbergzeta}), and the
\emph{Ruelle zeta} in the direct-product convention is $\Ruelle(s) =
\prod_\gamma(1-e^{-s\wl_\gamma})$.
\end{definition}

\begin{definition}[Guillemin flat trace]
\label{def:flat-trace}
For an Anosov flow with closed orbits $\gamma$, primitive periods
$\wl_\gamma$ and linearised Poincar\'e maps $\Poinc$, the \emph{flat trace}
of the Koopman operator is the distribution on $t>0$
$\flattr e^{-t\gflow} = \sum_\gamma\sum_{k\ge1}\wl_\gamma\,\delta(t-k\wl_\gamma)
/|\det(1-\Poinc^k)|$ (Dyatlov--Zworski, provenance
\texttt{prov:dz13-guillemin}). It is not a trace-class Hilbert-space trace.
\end{definition}

\begin{definition}[Flat dynamical determinant, first band]
\label{def:flat-tower}
$\lapfl(\lapvar) = \int_0^\infty e^{-\lapvar t}\,\flattr e^{-t\gflow}\,dt$ for
$\re\lapvar>0$, and the \emph{tower} $\Dtow(\lapvar) =
\prod_{j\ge1}\Zsel(\lapvar+j)$. \emph{Trace resonances} are the poles of the
continued $\lapfl$. The \emph{nonconstant first band} is
$\firstband = \{-\tfrac12\pm i\rj{j} : j\ge1\}$, with $\lamj{j} = \frac14 +
\rj{j}^2$ the Laplace eigenvalues and $\rj{0} = i/2$ for the constant mode
excluded.
\end{definition}

\begin{definition}[Scattering determinant of the modular group]
\label{def:scattering-determinant}
$\scatdet(s) = \sqrt\pi\,\Gamma(s-\tfrac12)\zeta(2s-1)/(\Gamma(s)\zeta(2s))$,
in the source's words \texttt{\detokenize{In the case when $\Gamma$ is the modular group $\mathrm{PSL}(2,\mathbb{Z})$, the scattering determinant is given by $$ \phi(s)=\sqrt{\pi}\frac{\Gamma(s-1/2)}{\Gamma(s)} \frac{\zeta(2s-1)}{\zeta(2s)},}} (provenance
\texttt{prov:fjs16-scattering-det}); the \emph{scattering multiplier} is
$\scmult(r) = -\tfrac12(\scatdet'/\scatdet)(\tfrac12+ir)$ and
$\cuspC = \mathcal F^{-1}\scmult$ with $\mathcal F^{-1}b(t) =
\frac1{2\pi}\int e^{irt}b(r)\,dr$.
\end{definition}

\begin{definition}[Prime comb measures]
\label{def:prime-measures}
$\primeP(t) = \sum_{n\ge2}\frac{\vM(n)}{n}\delta(t-2\log n)$ and
$\primech(t) = \sum_{n\ge2}\vM(n)n^{-1/2}\delta(t-2\log n)$ on $(0,\infty)$,
locally finite measures; $\archA$ is the archimedean-plus-boundary
distribution of \cref{prop:cusp-prime-comb}, equal to
$\tfrac12 - 1/(4\sinh(|t|/2))$ away from $t = 0$.
\end{definition}
